\section{CONTRIBUTIONS, LIMITATIONS AND FUTURE WORK}\label{ch:discussion}

\subsection{Contributions}\label{sec:disc-contrib}

The thesis makes three contributions. First, it evaluates regime-switching correlation for a stock--bond--gold portfolio with a joint VaR--ES loss, which the closest studies do not: \citet{paolella2019} use density and coverage scores on stocks, and \citet{fortin2023} use no regime switching. Second, it separates the volatility, correlation and mean channels under one regime signal. On the development data the gains come from modelling volatility, through regime-specific or GARCH volatility, and not from switching correlations, whose estimates were a third to a half of what the design could detect. Third, it reports a non-significant hold-out result with its interval, minimum detectable difference and power fixed in advance, so that a reader can see which effect sizes the data rule out and which they cannot.

Methodologically, the thesis shows that a common shortcut fails: the static test of regime-specific against pooled correlation is oversized when volatility clusters (Section~\ref{sec:corr-results}), so raw rejection counts across many specifications can look like evidence where there is none. Substantively, the hold-out shows that a desk choosing between these two correlation models is choosing between two forecasts that both under-forecast risk.

\subsection{What the evidence shows and does not show}\label{sec:disc-shows}

The pre-specified test did not reject (outcome O2). The data are consistent with a gain of the size the development data had suggested and with a zero difference, so the result is neither against regime-switching correlation nor for it. Three readings are not supported. That switching fails to improve forecasts: the one-sided lower bound, $-0.0033$, is below the smallest effect of interest of $-0.0020$. That switching helps: the test did not reject, and the 2022 estimate is too imprecise to say anything about the episode that motivated the question. And that the volatility regime, the correlation regime or the GARCH margins is responsible for a particular difference: several comparisons change more than one thing at once (Table~\ref{tab:differs}), and the non-nested ones compare model classes.

The development findings, that regime-specific volatility lowers the loss against a pooled normal, that GARCH margins do better still and that correlation steps are small, rest on many specifications examined after seeing results. I treat them as the reason the thesis has its shape, not as confirmed. The decomposition is the part I would defend most: for each step it gives an estimate, an interval and what the design could detect.

\subsection{What a risk desk should take from the hold-out}\label{sec:disc-calib}

The calibration result is the most useful for a practitioner, although it covers only the two primary models and its explanation is untested. Both breached the 5\% VaR on about 6.9\% of hold-out days, and on breach days losses exceeded the forecast expected shortfall. Which of the two is better is secondary to the fact that both fall short.

The pattern of the evidence points to the level of forecast volatility rather than to the correlation model or the tail shape. The portfolio's volatility rose from 7.1\% a year in the development period to 10.0\% in the hold-out (Table~\ref{tab:data-summary}), the Student-$t$ versions breached more often and not less, and the RiskMetrics EWMA, whose variance adapts within weeks, breached on 5.2\% of days with a loss that cannot be told apart from that of constant correlation. GARCH margins fitted on a five-year window and refitted every 63 days were slower to register the higher level of risk. This explanation fits the evidence but was not tested, and the EWMA comparison was chosen after the result was known. For a portfolio like this one, the evidence gives no reason to adopt regime-switching correlation, and checking the scale of the forecast against recent volatility matters more than the choice of correlation model.

\subsection{Reading the result against the economic framework}\label{sec:disc-slow}

The economic prediction was that, if the stock--bond correlation is driven by equity-market stress, a regime signal built from SPY volatility should help most when the recent correlation is positive. The hold-out does not bear this out. The conditional test did not reject; the favourable loss difference was larger on days with a negative lagged correlation; most of it accrued in 2021, when the calendar-year correlation was $-0.09$; and the switching model lost ground in the 2020 crash, the clearest equity-stress episode in the sample (Figure~\ref{fig:cumulative}). None of this is a test of the two readings, but it sits more easily with the macroeconomic account of \citet{campbell2020} and \citet{molenaar2024}, in which the sign of the correlation follows inflation and real rates and need not coincide with high equity volatility.

Two features of the design would slow any model of this kind. The stress-day correlation is learned from the estimation window, and before 2022 that window held almost no positive stock--bond correlation (Table~\ref{tab:eras}), so the model had little with which to represent a positive stress correlation however quickly the regime moved. And each model is estimated on the last 1,260 days and refitted every 21 days, so in 2022 the correlation changed within a window that still held several years of the old relationship. I did not test either explanation or run a faster-window or different-signal variant on the hold-out, because that analysis was not pre-specified.

\subsection{Limitations}\label{sec:disc-limits}

The following limits apply to every result in the thesis.

\begin{itemize}
\item \textbf{Power.} The test could detect a difference of 0.0031, which exceeds the smallest effect of interest of 0.0020, so outcome O3 required an estimate near or above zero (Section~\ref{sec:sesoi}). The test is better read as a measurement of how precisely this sample can speak than as a test with a realistic chance of excluding a meaningful gain. A longer ETF history would not add the regime of interest (Table~\ref{tab:eras}).
\item \textbf{One test, chosen with the development data in view.} The comparison was selected from a simulation in which the SPY--IEF shift is largest by construction, after I had examined ten universes and ten regime definitions. Nothing was fitted on the hold-out, but the specification is not independent of the exploration, and the pre-registration is a self-registration without an independent timestamp.
\item \textbf{A stylised power study.} The simulated process is an assumption, and simulated and real-data standard errors differ (Section~\ref{sec:power-results}), so the power figures describe the simulated process and need not transfer to markets. The 3.6\% rejection rate in the simulated null (500 replications) is a rough size check, and the upper-tail rate of 8.0\% at 1.645 shows that the normal reference is not exact there.
\item \textbf{Inference.} The estimand is the expected loss difference of two forecasting procedures, each with its own window, refit schedule and hidden-state chain, in the sense of \citet{giacomini2006}; it is not a statement about population parameters. The uncertainty of the estimated regimes is not carried into the test, since I did not bootstrap the whole pipeline. The 2020 and 2021 episodes dominate the average.
\item \textbf{Baselines.} The practitioner baselines, historical simulation and EWMA, were added after the hold-out was scored. They should have been part of the frozen design.
\item \textbf{Signal and window.} The regime signal is the SPY return only, and a volatility signal with a 1,260-day window may be slow to reflect a change in the stock--bond correlation.
\item \textbf{Scope.} There is one portfolio of three funds with constant weights, one data vendor and one reversal episode in the hold-out. The decision is made at the 5\% tail level, chosen for power, while regulatory backtests use the 1\% VaR. I make no claim about other portfolios, tail levels or markets.
\end{itemize}

\subsection{Future work}\label{sec:disc-future}

Each extension below would be a new study with its own pre-specification. A regime signal built from the stock--bond correlation itself or from inflation and real rates would test directly whether a slow macroeconomic state helps where a volatility state did not. A window that adapts to the regime, or exponentially weighted estimates, would address the slowness described above, and the EWMA result suggests that forecast scale should be part of any such comparison. A longer or broader sample, such as several stock--bond portfolios or international funds, would raise the power, and a power calculation like the one in Chapter~\ref{ch:eval} should come first. A bootstrap of the whole estimation pipeline would carry the uncertainty of the regime estimates into the test.
